\documentclass[letterpaper,journal]{IEEEtran}

\usepackage{amsmath,amsfonts}
\usepackage{array}
\usepackage{booktabs}
\usepackage{cite}
\usepackage{graphicx}
\usepackage{siunitx}
\usepackage{url}

\hyphenation{comag-ne-to-meter magneto-enceph-alo-graphy}
\sisetup{detect-all}

\begin{document}

\title{Measurement-Oriented Evaluation of Multilayer Thin-Ferrite Cylindrical Shields for Low-Noise Magnetic Instrumentation}

\author{Author~Name%
\thanks{Manuscript draft prepared for IEEE Transactions on Instrumentation and Measurement. Funding and author information will be added in the final manuscript.}
}

\markboth{IEEE Transactions on Instrumentation and Measurement}%
{Author: Multilayer Thin-Ferrite Cylindrical Shield Evaluation}

\maketitle

\begin{abstract}
Low-noise magnetic instruments are limited not only by sensor noise, but also by shield-induced leakage fields and shield-related magnetic noise. This manuscript develops a measurement-oriented finite-element evaluation framework for multilayer thin-ferrite cylindrical shields intended for ultra-sensitive magnetic measurement systems. The shield is treated as a device under test in a calibrated field configuration rather than as a passive enclosure. The framework separates two finite-element problems: an external uniform-field model for directional shielding-factor evaluation and a virtual pickup-coil model for the geometry-dependent magnetic-noise-related loss indicator $\int H^2 dV$ (IntH2). A validated no-shield directional reference is used to compute $\mathrm{SF}_x$, $\mathrm{SF}_z$, and the corresponding leakage ratios from field components rather than from vector magnitude. A compact analytical scaling argument is included to connect the finite-element matrix to thin cylindrical-shell theory without replacing the numerical model. The method also defines layer-resolved IntH2 decomposition, exact ferrite-volume calculation, volume-normalized auxiliary metrics, and a reproducible data-export and consistency-check pipeline. The current implementation targets manufacturable multilayer cylindrical thin-ferrite shells with \SI{100}{mm} inner radius, \SI{200}{mm} length, and ferrite sheet thickness from \SI{0.08}{mm} to \SI{0.60}{mm}. Validation studies required for IEEE TIM submission, including mesh and boundary-domain convergence, axial-field checks, segmented-shell correction, permeability sensitivity, and optional prototype shielding-factor measurement, are specified as part of the evaluation protocol. Numerical conclusions are reported only for cases with real AEDT exports and post-processing validation.
\end{abstract}

\begin{IEEEkeywords}
Ferrite magnetic shield, finite-element method, instrumentation, magnetic noise, multilayer cylindrical shell, shielding factor.
\end{IEEEkeywords}

\section{Introduction}
\IEEEPARstart{U}{ltra-sensitive} magnetic measurement systems, including atomic magnetometers, spin-exchange-relaxation-free (SERF) comagnetometers, magnetoencephalography instruments, and precision tests of fundamental symmetries, require a controlled low-field and low-noise magnetic environment \cite{ref1,ref2,ref3,ref4,ref5}. In such systems, the magnetic shield is part of the measurement chain. It determines the environmental-field leakage that reaches the sensor and can also generate magnetic-field noise through material loss. Therefore, the shield should not be evaluated only by a single shielding factor. For low-noise magnetic instrumentation, directional shielding factor, center leakage field, ferrite volume, radial space occupation, and material-loss-related noise indicators should be evaluated together.

High-permeability metallic shields are widely used to attenuate static and low-frequency fields \cite{ref6,ref7,ref8}. However, their high electrical conductivity produces eddy-current magnetic noise that can become relevant in the frequency band of high-sensitivity magnetometers \cite{ref9,ref10,ref11}. Ferrites are attractive alternatives because their conductivity is much lower than that of metallic alloys. This suppresses eddy-current noise; the remaining low-frequency shield-induced noise is then associated mainly with magnetic loss, represented by the imaginary permeability $\mu_r''$, and with the spatial distribution of magnetic field inside the ferrite \cite{ref9,ref12,ref13,ref14,ref15}. A recent IEEE TIM study on cubic ferrite shields demonstrated the practical importance of ferrites in compact low-noise shielding \cite{ref16}. Nevertheless, the cubic panel geometry does not directly address open-ended cylindrical shields, which are common around atomic sensors and compact comagnetometers.

A monolithic cylindrical ferrite shield is mechanically difficult to fabricate when the diameter is large and the wall is thin. Ferrites are brittle after sintering, and machining a \SI{200}{mm}-diameter thin-walled cylinder with sub-millimeter wall thickness is costly and fragile. A manufacturable alternative is to assemble a shield from multilayer thin-ferrite cylindrical shells or segmented ferrite sheets. This introduces a coupled design problem: layer number $N_L$, ferrite thickness $a$, radial gap $g$, circumferential segmentation, ferrite volume $V_f$, and total radial occupation all affect the center leakage field and the geometry-dependent loss integral used in magnetic-noise calculations.

Existing studies provide important foundations, but do not fully cover this measurement-oriented design problem. Kornack \textit{et al.} reported a low-noise ferrite cylinder with measured magnetic noise \cite{ref12}. Lee and Romalis derived magnetic-noise formulas for high-permeability shields and conducting objects with simple geometry \cite{ref9}. Later studies modeled cylindrical or multi-annular ferrite shields and discussed gap effects and material parameters \cite{ref13,ref14,ref17}. However, previous work has not simultaneously reported a validated no-shield directional reference, layer-resolved IntH2 analysis, and a multilayer thin-ferrite cylindrical geometry under a unified evaluation framework. Table~\ref{tab:comparison} summarizes the relationship between representative prior work and the present manuscript.

\begin{table*}[!t]
\caption{Relationship to Representative Ferrite Shield Studies}
\label{tab:comparison}
\centering
\footnotesize
\setlength{\tabcolsep}{3pt}
\begin{tabular}{p{0.16\textwidth}p{0.18\textwidth}p{0.18\textwidth}p{0.32\textwidth}}
\toprule
Work & Geometry & Main quantity & Scope relative to this work \\
\midrule
Kornack \textit{et al.} \cite{ref12} & Single ferrite cylinder & SF and measured noise & No multilayer thin-shell comparison \\
Lee and Romalis \cite{ref9} & Simple analytic shapes & Magnetic-noise formula & No finite-length multilayer FEM geometry \\
Ma \textit{et al.} \cite{ref13} & Cylindrical ferrite shield & SF and IntH2 & Layer-resolved IntH2 not reported \\
Lu \textit{et al.} \cite{ref14,ref17} & Multi-annular/gapped ferrite shields & Noise-related FEM analysis & Directional no-shield reference not emphasized \\
Sun \textit{et al.} \cite{ref16} & Cubic ferrite shield & SF and measured noise & Cubic geometry, not cylindrical multilayer shell \\
This work & Multilayer thin-ferrite cylinder & Directional SF, $V_f$, layerwise IntH2 & Measurement-oriented FEM and data pipeline \\
\bottomrule
\end{tabular}
\end{table*}

The novelty of this work is not simply the use of multiple ferrite layers. The contribution is a measurement-oriented evaluation framework that combines: 1) a validated no-shield directional reference for shielding-factor evaluation; 2) a virtual pickup-coil model for the geometry-dependent magnetic-noise-related loss integral; 3) layer-resolved IntH2 decomposition for multilayer shields; 4) fixed-volume comparison between single-layer and multilayer candidates; 5) unit-consistent auxiliary metrics for low-noise shield design; and 6) a reproducible raw-data, post-processing, validation, and plotting pipeline.

This draft focuses on the Introduction and Model/Method sections. Section~II defines the coordinate system, two finite-element models, directional shielding factor, analytical scaling motivation, virtual pickup-coil loss integral, cylindrical geometry, simulation matrix, data-export pipeline, auxiliary metrics, validation studies, and optional experimental protocol.

\section{Model and Evaluation Methodology}

\subsection{Coordinate System and Measurement Configuration}

The shield is modeled as a device under test (DUT) placed at the center of a calibrated magnetic-field region. The cylinder axis is defined as the Cartesian $z$-axis, while $x$ and $y$ span the transverse plane. Because the shield is a finite open-ended cylinder, transverse and axial shielding are not equivalent. The quantities $\mathrm{SF}_x$ and $\mathrm{SF}_y$ describe transverse shielding through the sidewall, whereas $\mathrm{SF}_z$ describes axial shielding affected by leakage through the open ends \cite{ref6,ref27,ref28,ref29}.

Two finite-element models are used, as shown schematically in Fig.~\ref{fig:framework}. These models answer different measurement questions and should not be physically conflated. The external uniform-field model evaluates environmental-field attenuation and provides $B_0$, $B_{\mathrm{center}}$, $\mathrm{SF}_q$, and $\mathrm{LeakageRatio}_q$. The virtual pickup-coil model evaluates how material loss couples to the measured magnetic field through the fluctuation-dissipation theorem and provides $\mathrm{IntH2}_{\mathrm{total}}$ and $\mathrm{IntH2}_{L_i}$.

\begin{figure}[!t]
\centering
\fbox{\parbox{0.92\columnwidth}{\centering\vspace{1.2cm}
Fig.~1 placeholder: measurement-oriented framework. External uniform-field model exports $B_0$ and $B_{\mathrm{center}}$ for SF calculation. Virtual pickup-coil model exports $\mathrm{IntH2}_{\mathrm{total}}$ and $\mathrm{IntH2}_{L_i}$. Post-processing computes $\eta_S$, $\eta_S^*$, $\rho_H$, and $\chi_H$.
\vspace{1.2cm}}}
\caption{Measurement-oriented evaluation workflow. The final manuscript will replace this placeholder with a compact vector diagram.}
\label{fig:framework}
\end{figure}

\subsection{Directional Shielding Factor}

For an external field applied along $q\in\{x,y,z\}$, the no-shield reference field at the observation point is
\begin{equation}
\mathbf{B}_0=(B_{0,x},B_{0,y},B_{0,z}),
\label{eq:B0}
\end{equation}
and the shielded center leakage field is
\begin{equation}
\mathbf{B}_{\mathrm{center}}=(B_{\mathrm{center},x},B_{\mathrm{center},y},B_{\mathrm{center},z}).
\label{eq:Bcenter}
\end{equation}
The directional shielding factor is defined by
\begin{equation}
\mathrm{SF}_q=\frac{|B_{0,q}|}{|B_{\mathrm{center},q}|},
\label{eq:sfq}
\end{equation}
and the corresponding leakage ratio is
\begin{equation}
\mathrm{LeakageRatio}_q=\frac{|B_{\mathrm{center},q}|}{|B_{0,q}|}=\frac{1}{\mathrm{SF}_q}.
\label{eq:leakage}
\end{equation}
The formal shielding factor is computed from the intended directional component. The vector magnitude $B_{\mathrm{Mag}}$ is retained only as an auxiliary consistency check and is not used as the primary SF definition.

\subsection{No-Shield Directional Reference}

The no-shield model is the calibration reference for the external-field problem. The ferrite bodies are deleted or assigned vacuum while the air region, boundary dimensions, observation point, solver setting, and excitation amplitude are kept identical to the shielded model. In vacuum, for $H_0=\SI{1}{A/m}$,
\begin{equation}
B_0=\mu_0 H_0=(4\pi\times10^{-7})(1)=1.256637\times10^{-6}\ \mathrm{T}.
\label{eq:vacuum}
\end{equation}
Thus, an ideal $x$-directed reference should give $B_{0,x}\approx\SI{1.256637e-6}{T}$, with $B_{0,y}$ and $B_{0,z}$ much smaller. The residual transverse components are treated as numerical residuals caused by finite-element discretization, finite-domain boundary approximation, and interpolation at the observation point.

The reference field is accepted only if the intended component dominates the other components:
\begin{equation}
\Gamma_q=\frac{|B_{0,q}|}{\max_{p\ne q}|B_{0,p}|}.
\label{eq:purity}
\end{equation}
In the processing pipeline, a failed directional reference prevents the corresponding shielding factors from being used in figures or conclusions.

\subsection{Analytical Scaling Motivation}

The finite-element model is the primary evaluation tool because the shield is finite in length, open ended, and may include layer gaps or segmented slots. Nevertheless, the choice of design variables is guided by the classical transverse cylindrical-shell solution. In a source-free region with linear isotropic permeability, $\nabla\times\mathbf{H}=0$ and $\nabla\cdot\mathbf{B}=0$ allow the magnetic scalar potential to satisfy Laplace's equation. For an infinitely long coaxial cylindrical shell in a transverse uniform field, the potential in each annular region can be written as
\begin{equation}
\phi_k(r,\theta)=\left(C_k r+\frac{D_k}{r}\right)\cos\theta ,
\label{eq:scalarpotential}
\end{equation}
where the constants are obtained from the continuity of normal $B$ and tangential $H$ at each interface \cite{ref6}. For a single high-permeability thin shell, this solution gives the familiar leading transverse scaling
\begin{equation}
S_T \simeq 1+\frac{\mu_r t}{2R},
\label{eq:singlelayerscaling}
\end{equation}
showing that shell thickness, mean radius, and permeability set the first-order shielding contribution.

For a multilayer coaxial shell, the analytical expansion contains direct layer terms and interlayer radius-coupling terms. A typical coupling factor between two layers with mean radii $R_i<R_j$ is
\begin{equation}
\Lambda_{ij}=1-\left(\frac{R_i}{R_j}\right)^2 .
\label{eq:couplingfactor}
\end{equation}
This motivates comparing layer number, radial gap, and total radial occupation rather than comparing total thickness alone. The present manuscript uses this analytical structure only as a scaling motivation; the reported $\mathrm{SF}_q$ values are obtained from the validated directional FEM reference in (\ref{eq:sfq}), not from the ideal infinite-cylinder expression.

Under the thin-shell approximation, a set of equal-height coaxial layers with total ferrite thickness $T_f=N_La$ has
\begin{equation}
V_f\approx 2\pi L R_c T_f ,
\label{eq:vfapprox}
\end{equation}
where $R_c$ is the radial-window center radius. Therefore, fixed $T_f$ is a useful first-order proxy for fixed ferrite volume. In the actual post-processing, however, $V_f$ is always computed from the exact layer radii in (\ref{eq:vf}) or from CAD volume export.

\subsection{Virtual Pickup-Coil Model and IntH2}

The shield-induced magnetic noise can be related to material loss through the fluctuation-dissipation theorem \cite{ref9,ref18}. A virtual pickup coil placed at the observation point produces a magnetic field $H_{\mathrm{vc}}$ in the ferrite. For coil area $A$, turn number $N_c$, peak current $I_0$, and magnetic moment $m=AN_cI_0$, the hysteresis-loss contribution to the magnetic-field noise amplitude spectral density can be written as
\begin{equation}
\delta B_{\mathrm{hyst}}(f)=
\frac{1}{m}
\left(
\frac{2k_{\mathrm{B}}T\mu''}{\pi f}
\int_V H_{\mathrm{vc}}^2\,dV
\right)^{1/2},
\label{eq:noisemodel}
\end{equation}
where $\mu''=\mu_0\mu_r''$ is the absolute imaginary permeability if $\mu_r''$ is specified as a relative quantity. The current $I_0$ is a peak current; if an RMS current is used in another implementation, the normalization must be converted consistently.

This manuscript does not claim measured magnetic-noise spectra. The exported FEM quantity is the geometry-dependent magnetic-noise-related loss indicator
\begin{equation}
\mathrm{IntH2}=\int_V H_{\mathrm{vc}}^2\,dV,
\label{eq:inth2}
\end{equation}
with SI unit $\mathrm{A^2\,m}$. IntH2 becomes an absolute noise prediction only after measured $\mu_r''(f)$, temperature, frequency, and virtual-coil normalization are specified. Relative comparison by IntH2 is meaningful only for identical material properties, identical coil normalization, and identical observation direction.

For two candidate geometries $A$ and $B$ evaluated under identical material loss and virtual-coil normalization, (\ref{eq:noisemodel}) gives the relative loss-induced field-noise scaling
\begin{equation}
\frac{\delta B_A(f)}{\delta B_B(f)}
=
\left(
\frac{\mathrm{IntH2}_A}{\mathrm{IntH2}_B}
\right)^{1/2}.
\label{eq:relnoise}
\end{equation}
Thus, IntH2 ratios should not be interpreted linearly as magnetic-noise ratios; the square-root relation applies only under the identical-normalization assumptions stated above.

For a multilayer shield, each ferrite layer is retained as an independent body. The layer-resolved terms are
\begin{equation}
\mathrm{IntH2}_{L_i}=\int_{V_{L_i}}H_{\mathrm{vc}}^2\,dV,
\label{eq:inth2layer}
\end{equation}
and
\begin{equation}
\mathrm{IntH2}_{\mathrm{total}}=\sum_{i=1}^{N_L}\mathrm{IntH2}_{L_i}.
\label{eq:inth2total}
\end{equation}
The export is accepted only if
\begin{equation}
\frac{|\mathrm{IntH2}_{\mathrm{total}}-\sum_i\mathrm{IntH2}_{L_i}|}
{|\mathrm{IntH2}_{\mathrm{total}}|}<10^{-3}.
\label{eq:inth2check}
\end{equation}

\subsection{Multilayer Cylindrical Geometry}

The DUT is an open-ended multilayer cylindrical shell with inner radius $R_{\mathrm{in}}$, axial length $L$, layer count $N_L$, ferrite thickness $a$, and radial interlayer gap $g$. The current design matrix uses
\begin{equation}
R_{\mathrm{in}}=\SI{100}{mm},\qquad L=\SI{200}{mm}.
\label{eq:basicgeometry}
\end{equation}
The total ferrite thickness and radial occupation are
\begin{equation}
T_f=N_La,
\label{eq:tf}
\end{equation}
\begin{equation}
T_{\mathrm{space}}=N_La+(N_L-1)g.
\label{eq:tspace}
\end{equation}
For layer $i$, where $i=1$ is the innermost layer,
\begin{align}
R_{i,\mathrm{in}} &= R_{\mathrm{in}}+(i-1)(a+g),\\
R_{i,\mathrm{out}} &= R_{i,\mathrm{in}}+a.
\end{align}
The ferrite volume is calculated from the exact cylindrical-shell expression
\begin{equation}
V_f=\pi L\sum_{i=1}^{N_L}\left(R_{i,\mathrm{out}}^2-R_{i,\mathrm{in}}^2\right).
\label{eq:vf}
\end{equation}
Although $V_f$ is displayed in $\mathrm{mm^3}$ for engineering readability, all SI-derived metrics are computed after unit conversion.

For segmented-shell correction, the circumferential coverage fraction is defined as
\begin{equation}
\eta_\phi=1-\frac{N_{\mathrm{seg}}g_\phi}{2\pi R_{\mathrm{mid}}},
\label{eq:etaphi}
\end{equation}
where $N_{\mathrm{seg}}$ is the number of circumferential segments, $g_\phi$ is the slot width, and $R_{\mathrm{mid}}$ is the layer mid-radius. Aligned and staggered slot patterns are evaluated separately when the segmented AEDT models are available.

\begin{figure}[!t]
\centering
\fbox{\parbox{0.92\columnwidth}{\centering\vspace{1.2cm}
Fig.~2 placeholder: multilayer cylindrical shell geometry with $R_{\mathrm{in}}$, $L$, $a$, $g$, $N_L$, $T_f$, and $T_{\mathrm{space}}$.
\vspace{1.2cm}}}
\caption{Geometry definition for the multilayer thin-ferrite cylindrical shell.}
\label{fig:geometry}
\end{figure}

\subsection{Simulation Matrix and Data Export}

The present main matrix is a transverse continuous-shell screening study. It is intentionally separated from axial, segmented, permeability-sensitivity, and convergence studies so that each validation stage can be tracked independently. Table~\ref{tab:matrix} lists the current main cases.

\begin{table*}[!t]
\caption{Main Continuous-Shell Transverse Matrix}
\label{tab:matrix}
\centering
\footnotesize
\setlength{\tabcolsep}{3pt}
\begin{tabular}{p{0.24\textwidth}ccccp{0.32\textwidth}}
\toprule
Case & $N_L$ & $a$ (mm) & $g$ (mm) & Direction & Purpose \\
\midrule
B0\_reference\_x & 0 & 0 & 0 & $x$ & no-shield reference \\
C1\_N1\_t008\_x & 1 & 0.08 & 0 & $x$ & thinnest single layer \\
C1\_N1\_t020\_x & 1 & 0.20 & 0 & $x$ & thin single layer \\
C1\_N1\_t040\_x & 1 & 0.40 & 0 & $x$ & intermediate single layer \\
C1\_N1\_t060\_x & 1 & 0.60 & 0 & $x$ & equal-volume single-layer baseline \\
C2\_N2\_t020\_g010\_x & 2 & 0.20 & 0.10 & $x$ & multilayer candidate \\
C2\_N3\_t020\_g010\_x & 3 & 0.20 & 0.10 & $x$ & equal-volume multilayer candidate \\
C2\_N4\_t015\_g008\_x & 4 & 0.15 & 0.08 & $x$ & equal-volume multilayer candidate \\
\bottomrule
\end{tabular}
\end{table*}

Each shielded external-field case exports $B_{\mathrm{center},x}$, $B_{\mathrm{center},y}$, $B_{\mathrm{center},z}$, and $B_{\mathrm{center,Mag}}$. The no-shield reference exports $B_{0,x}$, $B_{0,y}$, $B_{0,z}$, and $B_{0,\mathrm{Mag}}$. Each virtual pickup-coil case exports $\mathrm{IntH2}_{\mathrm{total}}$ and the layerwise quantities $\mathrm{IntH2}_{L_i}$. The post-processing scripts merge raw CSV exports by case ID and field direction, compute the derived metrics, and mark missing or failed rows without substituting zero values.

\subsection{Comparison Metrics}

The volume-normalized shielding index is defined as
\begin{equation}
\eta_S=\frac{\ln(\mathrm{SF}_x)}{V_f}.
\label{eq:etas}
\end{equation}
It is an auxiliary index introduced for this comparison, not a material property. When $V_f$ is reported in $\mathrm{mm^3}$, the displayed unit is $\mathrm{mm^{-3}}$; SI-derived quantities are computed after converting volume to $\mathrm{m^3}$.

A dimensionless normalized form is
\begin{equation}
\eta_S^*=
\frac{\ln(\mathrm{SF}_x)/V_f}{\ln(\mathrm{SF}_{x,\mathrm{ref}})/V_{\mathrm{ref}}},
\label{eq:etasstar}
\end{equation}
where the reference is the continuous single-layer shell with $a=\SI{0.60}{mm}$. The loss-integral density and attenuation-normalized loss indicator are
\begin{equation}
\rho_H=\frac{\mathrm{IntH2}_{\mathrm{total}}}{V_f[\mathrm{m^3}]},
\label{eq:rhoh}
\end{equation}
\begin{equation}
\chi_H=\frac{\mathrm{IntH2}_{\mathrm{total}}}{\ln(\mathrm{SF}_x)}.
\label{eq:chih}
\end{equation}
The metrics are interpreted jointly with $T_{\mathrm{space}}$ and the validation status. No single auxiliary metric defines a final shield geometry.

\subsection{Required Validation Studies}

Before the manuscript is extended to final results, the following validation studies are required.

\textit{Mesh convergence:} at least C1\_N1\_t008\_x and C2\_N4\_t015\_g008\_x are solved with coarse, medium, and fine mesh settings. The processed table records total element count, elements across ferrite thickness, $B_{\mathrm{center},x}$, SF$_x$, IntH2$_{\mathrm{total}}$, and relative changes. The target criteria are $\Delta\mathrm{SF}_x<1\%$, $\Delta\mathrm{IntH2}_{\mathrm{total}}<2\%$, and $\Delta B_{\mathrm{center},x}<1\%$ from medium to fine mesh.

\textit{Boundary-domain convergence:} C2\_N4\_t015\_g008\_x is evaluated with air-domain scales of $2R$, $3R$, and $5R$; optional $7R$ is retained as an additional check. The target criteria are less than 1\% change in SF$_x$ and less than 2\% change in IntH2$_{\mathrm{total}}$ from $3R$ to $5R$.

\textit{Axial shielding:} a validated B0\_reference\_z model is required before shielded SF$_z$ is computed. The selected equal-volume cases are C1\_N1\_t060\_z, C2\_N3\_t020\_g010\_z, and C2\_N4\_t015\_g008\_z. The transverse ranking is not generalized to axial shielding unless these cases support it.

\textit{Segmented-shell correction:} C2\_N4\_t015\_g008\_x is compared with segmented aligned and staggered slot models using $N_{\mathrm{seg}}=12$ and $g_\phi=\SI{1.6}{mm}$. The segmented-to-continuous penalty is quantified by SF$_{x,\mathrm{seg}}$/SF$_{x,\mathrm{cont}}$ and IntH2$_{\mathrm{seg}}$/IntH2$_{\mathrm{cont}}$.

\textit{Permeability sensitivity:} C1\_N1\_t060\_x and C2\_N4\_t015\_g008\_x are evaluated at $\mu_r'=500$, 1000, 2000, and 5000. The screening value $\mu_r'=1000$ is not treated as a measured material constant.

\textit{Field-map export:} the final paper should include external-field $B$ distributions, flux lines, $H_{\mathrm{vc}}$ contours, and $H_{\mathrm{vc}}^2$ contours to support the interpretation of flux redistribution and layerwise IntH2 contributions.

\textit{Prototype SF protocol:} if hardware is available, the prototype SF measurement uses a calibrated Helmholtz or triaxial coil system and a calibrated magnetic sensor. For independent $B_0$ and $B_{\mathrm{center}}$ measurements,
\begin{equation}
\frac{u(\mathrm{SF}_{q,\mathrm{exp}})}{\mathrm{SF}_{q,\mathrm{exp}}}
=
\sqrt{
\left[\frac{u(B_{0,q})}{B_{0,q}}\right]^2+
\left[\frac{u(B_{\mathrm{center},q})}{B_{\mathrm{center},q}}\right]^2
}.
\label{eq:uncertainty}
\end{equation}
For correlated measurements,
\begin{equation}
u_r^2(\mathrm{SF}_{q,\mathrm{exp}})=
u_r^2(B_0)+u_r^2(B_{\mathrm{center}})
-2\rho u_r(B_0)u_r(B_{\mathrm{center}}).
\label{eq:corruncertainty}
\end{equation}
If no experimental data are available, this section remains a protocol and is not presented as completed validation.

\section*{Acknowledgment}
Funding, facilities, and collaborator acknowledgments will be added in the final manuscript.

\begin{thebibliography}{31}

\bibitem{ref1}
J.~C.~Allred, R.~N.~Lyman, T.~W.~Kornack, and M.~V.~Romalis,
``High-sensitivity atomic magnetometer unaffected by spin-exchange relaxation,''
\textit{Phys. Rev. Lett.}, vol.~89, no.~13, Sep.~2002, Art.~no.~130801.

\bibitem{ref2}
I.~K.~Kominis, T.~W.~Kornack, J.~C.~Allred, and M.~V.~Romalis,
``A subfemtotesla multichannel atomic magnetometer,''
\textit{Nature}, vol.~422, no.~6932, pp.~596--599, Apr.~2003.

\bibitem{ref3}
M.~P.~Ledbetter, I.~M.~Savukov, V.~M.~Acosta, D.~Budker, and M.~V.~Romalis,
``Spin-exchange-relaxation-free magnetometry with Cs vapor,''
\textit{Phys. Rev. A}, vol.~77, no.~3, Mar.~2008, Art.~no.~033408.

\bibitem{ref4}
S.~Baillet,
``Magnetoencephalography for brain electrophysiology and imaging,''
\textit{Nature Neurosci.}, vol.~20, no.~3, pp.~327--339, Mar.~2017.

\bibitem{ref5}
T.~E.~Chupp, P.~Fierlinger, M.~J.~Ramsey-Musolf, and J.~T.~Singh,
``Electric dipole moments of atoms, molecules, nuclei, and particles,''
\textit{Rev. Mod. Phys.}, vol.~91, no.~1, Jan.~2019, Art.~no.~015001.

\bibitem{ref6}
T.~J.~Sumner, J.~M.~Pendlebury, and K.~F.~Smith,
``Conventional magnetic shielding,''
\textit{J. Phys. D: Appl. Phys.}, vol.~20, no.~9, pp.~1095--1101, Sep.~1987.

\bibitem{ref7}
V.~Kelha, J.~Pukki, R.~Peltonen, A.~Penttinen, R.~Ilmoniemi, and J.~Heino,
``Design, construction, and performance of a large-volume magnetic shield,''
\textit{IEEE Trans. Magn.}, vol.~18, no.~1, pp.~260--270, Jan.~1982.

\bibitem{ref8}
I.~Altarev \textit{et al.},
``A magnetically shielded room with ultra low residual field and gradient,''
\textit{Rev. Sci. Instrum.}, vol.~85, no.~7, Jul.~2014, Art.~no.~075106.

\bibitem{ref9}
S.-K.~Lee and M.~V.~Romalis,
``Calculation of magnetic field noise from high-permeability magnetic shields and conducting objects with simple geometry,''
\textit{J. Appl. Phys.}, vol.~103, no.~8, Apr.~2008, Art.~no.~084904.

\bibitem{ref10}
H.~B.~Dang, A.~C.~Maloof, and M.~V.~Romalis,
``Ultrahigh sensitivity magnetic field and magnetization measurements with an atomic magnetometer,''
\textit{Appl. Phys. Lett.}, vol.~97, no.~15, Oct.~2010, Art.~no.~151110.

\bibitem{ref11}
D.~Ma \textit{et al.},
``A novel low-noise mu-metal magnetic shield with winding shape,''
\textit{Sens. Actuators A, Phys.}, vol.~346, Oct.~2022, Art.~no.~113884.

\bibitem{ref12}
T.~W.~Kornack, S.~J.~Smullin, S.-K.~Lee, and M.~V.~Romalis,
``A low-noise ferrite magnetic shield,''
\textit{Appl. Phys. Lett.}, vol.~90, no.~22, May~2007, Art.~no.~223501.

\bibitem{ref13}
D.~Ma \textit{et al.},
``Parameter modeling analysis of a cylindrical ferrite magnetic shield to reduce magnetic noise,''
\textit{IEEE Trans. Ind. Electron.}, vol.~69, no.~1, pp.~991--998, Jan.~2022.

\bibitem{ref14}
J.~Lu \textit{et al.},
``Study of magnetic noise of a multi-annular ferrite shield,''
\textit{IEEE Access}, vol.~8, pp.~40918--40924, 2020.

\bibitem{ref15}
B.~Sun \textit{et al.},
``Correlating the microstructure of Mn--Zn ferrite with magnetic noise for magnetic shield applications,''
\textit{Ceram. Int.}, vol.~49, no.~8, pp.~11960--11967, Apr.~2023.

\bibitem{ref16}
B.~Sun, D.~Ma, X.~Fang, Y.~Xue, J.~Lu, H.~Chen, M.~Zhang, H.~Wei, B.~Han, and Y.~Zhai,
``Suppression of magnetic noise and field in cubic low-noise ferrite magnetic shields,''
\textit{IEEE Trans. Instrum. Meas.}, vol.~73, pp.~1--10, 2024, Art.~no.~1501810.

\bibitem{ref17}
J.~Lu \textit{et al.},
``Effect of gaps on magnetic noise of cylindrical ferrite shield,''
\textit{J. Phys. D: Appl. Phys.}, vol.~54, no.~25, Jun.~2021, Art.~no.~255002.

\bibitem{ref18}
R.~Kubo,
``The fluctuation-dissipation theorem,''
\textit{Rep. Prog. Phys.}, vol.~29, no.~1, pp.~255--284, Jan.~1966.

\bibitem{ref19}
C.~Liu \textit{et al.},
``Investigation on the effects of micro-vibration on the atomic comagnetometer,''
\textit{IEEE Trans. Instrum. Meas.}, vol.~72, 2023, Art.~no.~9511711.

\bibitem{ref20}
P.~Dutta and P.~M.~Horn,
``Low-frequency fluctuations in solids: $1/f$ noise,''
\textit{Rev. Mod. Phys.}, vol.~53, no.~3, pp.~497--516, Jul.~1981.

\bibitem{ref21}
K.~Yang \textit{et al.},
``Improved measurement of the low-frequency complex permeability of ferrite annulus for low-noise magnetic shielding,''
\textit{IEEE Access}, vol.~7, pp.~126059--126065, 2019.

\bibitem{ref22}
K.~Yang \textit{et al.},
``Minimizing magnetic fields of the low-noise MnZn ferrite magnetic shield for atomic magnetometer,''
\textit{J. Phys. D: Appl. Phys.}, vol.~55, no.~1, Jan.~2022, Art.~no.~015003.

\bibitem{ref23}
Z.~Xu, Z.~Zhang, B.~Wu, H.~Wang, X.~Kong, and M.~Wang,
``A study of enclosed magnetic shielding room by simulation,''
\textit{IEEE Trans. Appl. Supercond.}, vol.~31, no.~8, pp.~1--5, Nov.~2021.

\bibitem{ref24}
X.~Xu, L.~Wang, W.~Liu, and Z.~Zhao,
``Theoretical modeling and characterization of equivalent magnetic properties in laminated composite magnetic shielding,''
\textit{Measurement}, vol.~251, 2025, Art.~no.~117237.

\bibitem{ref25}
\textit{Cores Made of Soft Magnetic Materials---Measuring Methods---Part 2: Magnetic Properties at Low Excitation Level}, IEC 62044-2, 2005.

\bibitem{ref26}
TDK Corporation,
``Mn--Zn ferrite material characteristics: PC95 and related Mn--Zn ferrite grades,''
product catalog, accessed May~2026. [Online]. Available: \url{https://product.tdk.com/system/files/dam/doc/product/ferrite/ferrite/ferrite-core/catalog/ferrite_mn-zn_material_characteristics_en.pdf}

\bibitem{ref27}
A.~Mager,
``Magnetic shields,''
\textit{IEEE Trans. Magn.}, vol.~6, no.~1, pp.~67--75, Mar.~1970.

\bibitem{ref28}
K.~Nagashima, I.~Sasada, and K.~Tashiro,
``High-performance bench-top cylindrical magnetic shield with magnetic shaking enhancement,''
\textit{IEEE Trans. Magn.}, vol.~38, no.~5, pp.~3335--3337, Sep.~2002.

\bibitem{ref29}
J.~M.~G.~D.~K.~O'Connell, J.~M.~Brown, and T.~W.~Kornack,
``Technique for high axial shielding factor performance of large-scale, thin, open-ended, cylindrical Metglas magnetic shields,''
arXiv:1107.2625, 2011.

\bibitem{ref30}
J.~Fang, S.~Wan, J.~Qin, and Y.~Chen,
``A novel method for in-situ measuring the static shielding factor of a magnetic shield,''
\textit{Measurement}, vol.~170, 2021, Art.~no.~108718.

\bibitem{ref31}
T.~Brys \textit{et al.},
``Magnetic field stabilization for magnetically shielded volumes by external field coils,''
\textit{J. Appl. Phys.}, vol.~116, no.~8, 2014, Art.~no.~084903.

\end{thebibliography}

\end{document}
